\subsection{Split complexes}

PROPOSITION 6. — Let (C, d) be a complex. The following conditions are equivalent:

\begin{enumerate}
\def\labelenumi{(\roman{enumi})}
\item
  there exists a homotopism from (C, d) to (H(C), 0)
\item
  there exists a graded A-endomorphism s of C, of degree 1, such that \(d = d \circ s \circ d\) ;
\item
  B(C) and Z(C) are direct summands of C;
\item
  (C, d) is a direct sum of subcomplexes that are either of length 0, or of length 1 with zero homology.
\item
  \(\Rightarrow\) (ii): let \(\varphi: \mathbf{C} \to \mathbf{H}(\mathbf{C})\) be a homotopism; there then exists a morphism of complexes \(\psi: \mathbf{H}(\mathbf{C}) \to \mathbf{C}\) and an endomorphism \(s\) of \(\mathbf{C}\) of degree 1 such that \(\psi \circ \varphi = 1_{\mathbf{C}} - s \circ d - d \circ s\) . We have \(d \circ \psi = \psi \circ 0 = 0\) , hence
\end{enumerate}

\[
0 = d \circ \psi \circ \varphi = d - d \circ s \circ d - d \circ d \circ s = d - d \circ s \circ d, \quad \text {whence (ii)}.
\]

\begin{enumerate}
\def\labelenumi{(\roman{enumi})}
\setcounter{enumi}{1}
\item
  \(\Rightarrow\) (iii): let \(s\) be as in (ii). Then \(d \circ (1_{\mathbb{C}} - s \circ d) = 0\) , hence \(1_{\mathbb{C}} - s \circ d\) is a projector of \(\mathbf{C}\) onto \(\mathbf{Z}(\mathbf{C})\) , and \((d \circ s) \circ d = d\) , hence \(d \circ s\) is a projector of \(\mathbf{C}\) onto \(\mathbf{B}(\mathbf{C})\) .
\item
  \(\Rightarrow\) (iv): for each \(n\in\mathbf{Z}\) , set \(Z_{n}=Z_{n}(\mathrm{C})\) , \(B_{n}=B_{n}(\mathrm{C})\) and choose submodules \(K_{n}\) and \(B_{n}^{\prime}\) of \(C_{n}\) such that \(C_{n}=B_{n}^{\prime}\oplus Z_{n}\) , \(Z_{n}=K_{n}\oplus B_{n}\) . Then
\end{enumerate}

\[
\mathrm{E} _ {(n)} = \mathrm{K} _ {n} (- n) \quad \text { and } \quad \mathrm{F} _ {(n)} = \mathrm{B} _ {n} ^ {\prime} (- n) \oplus \mathrm{B} _ {n - 1} (1 - n)
\]

are subcomplexes of (C, d); we have (C, d) = \(\bigoplus_{n \in \mathbf{Z}} (\mathrm{E}_{(n)} \oplus \mathrm{F}_{(n)})\) ; each \(\mathrm{E}_{(n)}\) is either zero or of length 0, each \(\mathrm{F}_{(n)}\) is either zero or of length 1 and with zero homology, whence (iv).

\begin{enumerate}
\def\labelenumi{(\roman{enumi})}
\setcounter{enumi}{3}
\tightlist
\item
  \(\Rightarrow\) (i): it suffices to note that (i) is satisfied when C is of length zero, or of zero homology and length 1.
\end{enumerate}

DEFINITION 6. --- A complex C is said to be split if it satisfies the equivalent conditions of prop. 6.

An endomorphism s of C satisfying condition (ii) of prop. 6 is called a splitting of C.

Examples. --- 1) A complex with zero differential is split.

\begin{enumerate}
\def\labelenumi{\arabic{enumi})}
\setcounter{enumi}{1}
\item
  The complexes homotopic to zero are the split complexes with zero homology, i.e.~the complexes C such that H(C) = 0 and Z(C) is a direct summand of C.
\item
  Let \(f: \mathbf{M} \to \mathbf{N}\) be a homomorphism of A-modules and C the complex such that \(\mathbf{C}_1 = \mathbf{M}, \mathbf{C}_0 = \mathbf{N}, \mathbf{C}_i = 0\) for \(i \neq 0, 1, d_1 = f, d_i = 0\) for \(i \neq 1\) . Then C is split if and only if Ker \(f\) is a direct summand of M and Im \(f\) a direct summand of N.
\item
  The complex C is split as soon as B(C) and H(C) are projective (resp. as soon as \(B_{n}(C)\) and \(H_{n}(C)\) are injective for each n). Indeed, by the exact sequences \((I_{n})\) to \((IV_{n})\) of no. 1, Z(C) is then a direct summand of C and B(C) a direct summand of Z(C) (resp. B(C) is a direct summand of C and Z(C)/B(C) a direct summand of C/B(C)).
\item
  In particular, if A is principal, a free complex C is split if and only if H(C) is free (that is, H \(_{n}\) (C) free for every n ∈ Z).
\end{enumerate}

Remark. --- a) Suppose that the canonical exact sequence of graded A-modules

\[
0 \to \mathrm{B} (\mathrm{C}) \to \mathrm{Z} (\mathrm{C}) \xrightarrow {\pi} \mathrm{H} (\mathrm{C}) \to 0\tag{II}
\]

is split (this holds for example if H(C) is projective, or B \(_{n}\) (C) injective for each n); let σ : H(C) → Z(C) be a graded A-linear section of π, and let ψ be the homomorphism x ↦ σ(x) from H(C) to C. Then ψ is a homologism from (H(C), 0) to C, such that H(ψ) = 1 \(_{H(C)}\) .

\begin{enumerate}
\def\labelenumi{\alph{enumi})}
\setcounter{enumi}{1}
\tightlist
\item
  Suppose that the canonical exact sequence of graded A-modules
\end{enumerate}

\[
0 \rightarrow \mathrm{H} (\mathrm{C}) \stackrel {i} {\rightarrow} \mathrm{C} / \mathrm{B} (\mathrm{C}) \rightarrow \mathrm{B} (\mathrm{C}) (- 1) \rightarrow 0\tag{IV}
\]

is split (this occurs, for example, if B(C) is projective, or H \(_{n}\) (C) is injective for each \(n\) ); let \(\tau: \mathbf{C} / \mathbf{B}(\mathbf{C}) \to \mathbf{H}(\mathbf{C})\) be a graded A-linear retraction of \(i\) and \(\varphi\) the homomorphism from C to \(\mathbf{H}(\mathbf{C})\) which associates to each element of C the image under \(\tau\) of its class modulo \(\mathbf{B}(\mathbf{C})\) . Then \(\varphi\) is a homologism from C to \((\mathbf{H}(\mathbf{C}), 0)\) such that \(\mathbf{H}(\varphi) = 1_{\mathbf{H}(\mathbf{C})}\) .

